\section*{Results}

\subsection*{What parity guarantees, and what rotation alone still guarantees}

Group theory fixes, before any measurement, what each symmetry class can promise (proofs in \snote{snote:proofs}).

A periodic structure $x$ carries $n$ atoms at positions $\mathbf{r}_1,\dots,\mathbf{r}_n$ with species labels, on a Bravais lattice $\Lambda$. A model $f_\theta$ with parameters $\theta$ maps $x$ into the target space $V$. For a rank-$r$ polar tensor, $V$ carries the representation $D(Q) = Q^{\otimes r}$ of $Q \in \mathrm{O}(3)$, so that $D(I) = (-1)^{r}\,\mathrm{id}$ at the inversion $I$. The piezoelectric tensor has $r = 3$ and $V = \mathbb{R}^{3}\otimes\mathrm{Sym}^{2}(\mathbb{R}^{3})$, the eighteen-dimensional space of rank-3 tensors symmetric in the strain indices. We write $\lVert\cdot\rVert_F$ for the Frobenius norm, so that $\lVert f_\theta(x)\rVert_F$ is the violation magnitude $\lVert e \rVert_F$ that every experiment below reports.

Three assumptions are used throughout. \textbf{(A1)} \emph{Permutation invariance}: relabelling atoms of the same species leaves the output unchanged. \textbf{(A2)} \emph{Translation invariance}: a rigid shift by any vector, and a shift of any atom by a lattice vector, leave it unchanged. \textbf{(A3)} \emph{$G$-equivariance}: $f_\theta(Q\!\cdot\!x)=D(Q)f_\theta(x)$ for all $Q\in G$, with $G=\mathrm{O}(3)$ for a parity-labelled network and $G=\mathrm{SO}(3)$ otherwise. The rigid-shift clause of (A2) is load-bearing, because a non-symmorphic space group pairs a point-group rotation with a fractional translation that is not a lattice vector. Every architecture here satisfies (A1) and (A2) by construction, verified numerically.

\begin{definition}[Centrosymmetric structure]\label{def:centro}
$x$ is centrosymmetric if there is a point about which the inversion maps the structure onto itself. Placing that point at the origin, this means there is a species-preserving permutation $\pi$ of the atoms and lattice vectors $\mathbf{t}_i \in \Lambda$ with $-\mathbf{r}_i = \mathbf{r}_{\pi(i)} + \mathbf{t}_i$ for every $i$. Because $I^2 = \mathrm{id}$, the permutation is an involution, $\pi^{2} = \mathrm{id}$.
\end{definition}

\begin{lemma}[Inversion identity]\label{lem:inv}
Let $f_\theta$ satisfy (A1) and (A2). Then $f_\theta(I\!\cdot\!x) = f_\theta(x)$ for every centrosymmetric $x$ and every $\theta$.
\end{lemma}

The inverted copy of a centrosymmetric crystal is the same periodic solid, relabelled and shifted, so a model that cannot see relabelling or shifts cannot see the inversion. By (A2) the choice of inversion centre is immaterial. Adding parity to this identity gives the guarantee.

\begin{theorem}[Structural zero]\label{thm:zero}
Let $f_\theta$ satisfy (A1), (A2) and (A3) with $G = \mathrm{O}(3)$, and let the target be a rank-$r$ polar tensor with $r$ odd. Then $f_\theta(x) = 0$ for every centrosymmetric $x$ and every $\theta$.
\end{theorem}

The proof mentions neither the parameters, the data nor the loss, which is what makes the zero structural rather than learned; for an SO(3)-labelled network its second step is unavailable, because $D(I)$ is never invoked. In exact arithmetic that is the whole story; in floating point, the residual is bounded.

Three consequences follow, stated formally and proved in the Supplementary Information. The residual is bounded by the measured equivariance defects, $\lVert f_\theta(x)\rVert_F \le (\delta + \eta)/2$ (Proposition~\ref{sprop:floor}); for an O(3) model both defects are set by the arithmetic, so its floor tracks precision and not symmetry, which is why the measured floors sit at the single-precision level, while for a parity-blind model the bound is a tautology. Test-time antisymmetrization, $T_{\mathrm{sym}}(x) = [f_\theta(x) - f_\theta(I\!\cdot\!x)]/2$, vanishes identically on every centrosymmetric $x$ whether or not the features are parity-labelled, by Lemma~\ref{lem:inv} alone, so the zero carries no parity information and degrades in proportion to $\delta$. And the guarantee reaches derivatives: the Jacobian satisfies $J\circ P = -J$ for the inversion map $P$ on displacement fields, hence $J P_{\mathrm{even}} = 0$, so the even-subspace energy fraction vanishes.

\begin{corollary}[Rotation ceiling]\label{cor:rot}
Let $f_\theta$ satisfy (A1)--(A2) and (A3) with $G=\mathrm{SO}(3)$ only, and let $G^{+}(x)\subset\mathrm{SO}(3)$ be the proper-rotation subgroup of the point group of $x$; each of its elements is the rotational part of a space-group operation of $x$, and so returns the crystal to itself once the accompanying translation, which may be fractional, is absorbed by (A2). Then $f_\theta(x)$ lies in the invariant subspace $V^{G^{+}(x)} = \{T \in V : R^{\otimes r}T = T\ \forall R \in G^{+}(x)\}$. If $V^{G^{+}(x)} = \{0\}$, then $f_\theta(x)=0$ is forced with no parity information; otherwise nothing forces it to vanish. If (A1)--(A3) hold only approximately over $G^{+}(x)$, the forced case weakens to a bound proportional to the equivariance defect.
\end{corollary}

For the piezoelectric tensor this is a sharp, parameter-free prediction. Among the centrosymmetric classes, only m$\bar 3$m has a proper-rotation subgroup, 432, that annihilates every rank-3 piezoelectric tensor, so Corollary~\ref{cor:rot} forces a rotation-only model to zero on m$\bar 3$m crystals and no others: 166 of the 2,000 evaluation crystals, leaving 1,834 (91.7\%) on which only parity can enforce the zero. The ceiling is exact only for an exactly equivariant model.

Theorem~\ref{thm:zero} says an O(3) model \emph{must} predict zero on a centrosymmetric crystal; the rotation ceiling says an SO(3) model \emph{may} predict a nonzero tensor wherever rotation alone does not forbid one. That it does, on almost every crystal where it is free to, is empirical.

\subsection*{Parity labels decide whether impossible predictions occur}

Four targets separate two properties that a single experiment would confound: being a parity-odd target, and being evaluated where symmetry forbids a nonzero value. The QM9~\cite{ramakrishnan2014qm9} internal energy $U_0$ (parity-even scalar) and the elasticity tensor (parity-even rank 4) carry no odd-parity constraint. The QM9 dipole is parity-odd, and a point-group census of the evaluation population finds that symmetry forces the dipole to vanish for only 7 of its 10,831 molecules (0.065\%), methane among them, so the constraint is inactive on all but a negligible fraction (\snote{snote:qm9pg}). Only the piezoelectric tensor on centrosymmetric crystals is both (Fig.~\ref{fig:concept}a). The evaluation population is 2,000 centrosymmetric insulators from the Materials Project~\cite{jain2013materials}, re-verified with spglib~\cite{togo2024spglib}, in two coordinate variants: idealized (snapped onto the exact space group) and raw (as relaxed by density functional theory).

The measurement is a matched-pair ablation (Fig.~\ref{fig:concept}b). NequIP, Allegro and MACE were each built in two arms through a single irreducible-representation helper: an O(3) arm carrying the natural spherical-harmonic parity $(-1)^{l}$, and an SO(3) arm in which every feature is declared parity-even. Only the parity bookkeeping differs, and with it the reflection law the feature algebra enforces; architecture, hyperparameters, data, splits and seeds are shared. EqV2 joins unmodified as a deployed rotation-only representative, with no parity labels to toggle. The grid is 7 arms $\times$ 4 targets $\times$ 3 seeds totaling 84 runs.

Before training, every arm passed a verification gate under a random rotation and a random improper operation: O(3) arms must satisfy both transformation laws, SO(3) arms must satisfy rotation and fail reflection. All did, and the SO(3) parameter count exceeds its O(3) counterpart in every configuration (ratios 1.00--1.40), so reduced capacity cannot explain any SO(3) deficit (per-configuration counts accompany the code release).

\begin{figure}[t]
\centering
\includegraphics[width=\textwidth]{fig1_concept}
\caption{A label-free test of physical validity. a, A centrosymmetric crystal, shown with its unit cell (dashed). Every atom has a partner at $-\mathbf{r}$ of the same species, reached by the arrows through the crossed circle at the inversion centre, so the inversion maps the crystal onto itself, $I\!\cdot\!x = x$. The piezoelectric tensor is parity-odd, $e \mapsto -e$, and Neumann's principle requires it to be invariant under every operation of the point group; the two demands together force $e = -e$, hence $e = 0$ exactly. b, The two networks are built from one code path and differ in one bit: the feature irreducible representations are $0e,\,1o,\,2e,\,3o$ up to $l_{\max}=3$ in the O(3) arm, and every degree even in the SO(3) arm. c, Predicted magnitudes on the 2,000 spglib-verified centrosymmetric crystals, for the NequIP pair; boxes span the 25th to 75th percentile, whiskers the 5th to 95th, and the horizontal line is the median. The true value is exactly zero for every structure. The O(3) median sits at the arithmetic floor, $3\times10^{-7}$\,C\,m$^{-2}$; the SO(3) arm predicts tensors above the operating threshold.}
\label{fig:concept}
\end{figure}

The two arms separate by six orders of magnitude on the idealized evaluation population, where the true value is exactly zero (Fig.~\ref{fig:concept}c for the NequIP pair, all cores in Fig.~\ref{fig:headline}a). Not one of the 18,000 O(3) structure-runs exceeds the operating threshold of 0.01\,C\,m$^{-2}$: their medians sit at the arithmetic floor, and the largest single prediction falls a factor of seven below the threshold. The SO(3) arms cross it for 90--96\% of the population, predictions we call false flags, and their median violation matches the magnitude of the genuine piezoelectric tensors in the training data: the typical SO(3) prediction is the size of a real response, not a numerical residue. Per-core fractions, medians and the threshold calibration are tabulated in the Supplementary Information. Because the summed readout is extensive while the piezoelectric tensor is intensive, the whole grid was also trained with a mean-pooled readout and a threshold recalibrated on the mean-pooled distribution. The O(3) arms return a false-flag fraction of exactly zero under both poolings on every core and seed, and the SO(3) fractions move by a mean of $0.0065$ across twelve matched pairs, less than the seed-to-seed spread within either arm (\snote{snote:pooling}): neither the guarantee nor its absence depends on how the readout is normalized.

The separation depends on neither the operating point nor the coordinate variant. Across four decades of threshold the two arms stay cleanly separated (Fig.~\ref{fig:headline}b), and the SO(3) false-flag fractions are unchanged between idealized and raw coordinates. The one nuance is on the O(3) side: a single structure produces a small raw-variant response and proves genuinely non-centrosymmetric below the selection tolerance, so the guarantee detected a data artefact.

The separation is statistically decisive: structure-paired Wilcoxon tests reject equality at every core and seed, per-structure violations correlate across seeds, and independent architectures flag almost the same crystals, so the failure is a property of the symmetry class rather than of an implementation.

\begin{figure}[t]
\centering
\includegraphics[width=\textwidth]{fig2_headline}
\caption{Parity labels separate possible from impossible predictions by six orders of magnitude. a, Distribution of the predicted piezoelectric magnitude $\lVert e \rVert_F$ on 2,000 spglib-verified centrosymmetric crystals (idealized coordinates), for the O(3) and SO(3) arms of each core. Symmetry forces the true value to zero, so the shaded region above the operating threshold $\tau$ marks physically impossible predictions. Boxes span the 25th to 75th percentile of the per-structure distribution (mean of per-seed percentiles over 3 seeds), whiskers the 5th to 95th, the horizontal line is the median, and the isolated dash is the maximum. Percentages above each SO(3) box give the fraction of structures exceeding $\tau$. The double arrow gives the ratio of the two NequIP medians, $2\times10^{6}$. b, False-flag fraction as a function of threshold (idealized coordinates), mean over 3 seeds; shaded bands show one standard deviation. Blue, O(3) arms; orange, SO(3) arms; markers distinguish cores. EqV2 has no O(3) arm.}
\label{fig:headline}
\end{figure}

\subsection*{The constraint costs nothing where it is inactive}

If parity labelling handicapped models on ordinary tasks, removing it could be a defensible trade. It does not. On the three targets where no symmetry-forbidden value is evaluated ($U_0$, dipole and elasticity), the two arms are indistinguishable within seed noise, and any visible tendency favours keeping the labels: on the dipole and on elasticity the O(3) arm is nominally better in all three cores, by 0.04 to 0.74 pooled standard deviations, an effect three seeds cannot resolve and from which we claim no significance, while on $U_0$ the sign varies across cores (Fig.~\ref{fig:accuracy}a). The dipole is the decisive control: it is parity-odd like the piezoelectric tensor, yet the arms match, because the constraint is inactive for all but a negligible fraction of the test population, far too few to move an aggregate error. What separates the arms is not the target's parity but evaluation where symmetry forbids a nonzero value.

Where the constraint is active, it helps rather than costs: on the piezoelectric tensor the O(3) arm has the lower in-distribution error in all three matched cores, with a consistent sign (Fig.~\ref{fig:accuracy}b), and does so while carrying fewer parameters than its SO(3) counterpart in every configuration. EqV2's in-distribution error is comparable, so its near-universal false flagging coexists with ordinary competence on non-centrosymmetric crystals. Accuracy tables, effect sizes and the compute budget are in \snote{snote:accuracy}.

\begin{figure}[t]
\centering
\includegraphics[width=\textwidth]{fig3_accuracy}
\caption{Removing parity labels costs nothing where the parity constraint is inactive, and costs accuracy where it is active. a, Difference in test mean absolute error (MAE) between the SO(3) and O(3) arms, in units of the pooled per-seed standard deviation, for the three matched cores on all four targets. Positive values favour the O(3) arm. The grey band marks one pooled standard deviation, within which a three-seed comparison cannot resolve a difference; the tinted column is the one target on which the parity constraint is active. b, Piezoelectric test MAE for both arms of each core (mean over 3 seeds; error bars, one standard deviation), with the two arms of a core joined and the effect size printed above. EqV2 has no O(3) arm.}
\label{fig:accuracy}
\end{figure}

\subsection*{Training cannot buy the zero}

The natural objection is that the SO(3) failure is a data problem: no centrosymmetric crystal appears in the piezoelectric training set. We retrained every SO(3) arm on the training set augmented with 1,000 spglib-verified centrosymmetric insulators labelled with the exact zero tensor, drawn from six space groups, with hyperparameters and partitions unchanged.

The decisive measurement is the in-training control: predictions on the very crystals whose explicit zero labels the model was trained on. After training, the models still false-flag roughly nine in ten of them, 0.895--0.922 across the four cores. This is not underfitting, because the zero-labelled rows are fit better than the real-tensor rows in training error; nor is it a coverage problem, since held-out false-flag fractions barely move whether or not the crystal's space group was seen in training. Gradient descent drives the predictions towards zero and cannot arrive, because the required value is exact; augmentation meanwhile improves the regression error for every core, producing better regressors that predict impossible values.

Scale does not close the gap either. In a zero-injection sweep of the NequIP SO(3) arm from 250 to 16,000 zero-labelled crystals, the held-out violation median falls 8.7-fold, from 0.666 to 0.077\,C\,m$^{-2}$, while the held-out false-flag fraction moves only from 0.895 to 0.858 (three seeds per size): the predictions shrink, but almost none becomes zero. Weighting those rows a hundredfold in the loss is no better: it reduces the trained-on false-flag fraction by 0.22, leaving two-thirds flagged, and never approaches the exact zero the O(3) arms deliver at any weight, for free. All three experiments, with per-epoch training curves, are tabulated in \snote{snote:augmentation}.

A second repair operates at test time: antisymmetrize the output over inversion, $T_{\mathrm{sym}}(x) = [T(x) - T(I\!\cdot\!x)]/2$. On the exactly-equivariant SO(3) cores this drives the false-flag fraction to zero, and the reason the success is empty was proved above: by the inversion identity the construction vanishes identically for \emph{any} permutation- and translation-invariant model, so no parity information is involved. It also presupposes knowing the target's parity in advance, precisely the knowledge O(3) features carry, and it fails on EqV2, whose approximate equivariance violates the permutation identity the fix relies on: its false-flag fraction falls only to 0.82.

\subsection*{The guarantee tracks the physics of symmetry breaking}

First, the guarantee switches off when the symmetry does. We displaced rutile TiO$_2$ along a polar mode over 33 symmetry-verified amplitudes, from the centrosymmetric parent to beyond the physical distortion. The O(3) arms rise smoothly from the arithmetic floor, while the SO(3) arms sit on a spurious floor five to six orders of magnitude higher, already present at zero distortion and dominating the physical signal throughout the small-distortion regime. Rutile is decisive because its rotation subgroup 422 is piezoelectric, so by the rotation ceiling only parity forbids a response there; the separation reproduces on two further 4/mmm polar paths, rutile-type SnO$_2$ and anatase TiO$_2$.

Second, the crystals SO(3) models get right are exactly those the rotation ceiling identifies (Fig.~\ref{fig:mechanism}a). The three matched cores flag none of the m$\bar 3$m crystals, whose rotation subgroup alone forbids the tensor, and almost all of those on which nothing but parity can enforce the zero: an exactly rotation-equivariant model saturates the ceiling. EqV2 exceeds it, as the ceiling's stability clause allows for approximate equivariance, and the excess is entirely its spurious flagging of m$\bar 3$m crystals an exact model cannot flag at all, 53.4\% of them. Per-family fractions are in \snote{snote:named}.

Third, the Jacobian identity above requires the even-subspace energy fraction of the Jacobian to be exactly zero at a centrosymmetric geometry. Over 20 crystals spanning 19 space groups, with every Jacobian certified against finite differences, the two arms again sit five to six orders of magnitude apart (Fig.~\ref{fig:mechanism}b): the O(3) arm at the arithmetic floor, the SO(3) arm at order unity. The O(3) architecture does not merely output a zero; it encodes which atomic displacements can activate a piezoelectric response.

The same pattern shows on named materials (Fig.~\ref{fig:mechanism}c). The SO(3) arms predict a full-magnitude tensor for corundum Al$_2$O$_3$, whose non-cubic point group leaves parity as the only protection, and near zero for the rocksalt and diamond-structure compounds of class m$\bar 3$m. On those m$\bar 3$m rows the 72 matched-core SO(3) predictions (3 cores $\times$ 3 seeds on eight compounds) span $5.7\times10^{-10}$ to $3.2\times10^{-6}$\,C\,m$^{-2}$. EqV2 exceeds the threshold on two of them, diamond and SrTiO$_3$, the approximate-equivariance failure recorded by its m$\bar 3$m bar in Fig.~\ref{fig:mechanism}a.

\begin{figure}[t]
\centering
\includegraphics[width=\textwidth]{fig5_mechanism}
\caption{The missing guarantee is specifically parity, and it is differentiable. a, False-flag fraction of the SO(3) arms split by the point-group family of the evaluated crystal: m$\bar 3$m, whose proper-rotation subgroup 432 alone forbids a rank-3 piezoelectric tensor; m$\bar 3$, whose subgroup 23 does not; and all non-cubic classes, none of whose rotation subgroups forbid it. The O(3) arms are zero in every family. b, Even-subspace energy fraction of the Jacobian $J = \partial T/\partial \mathbf{r}$ at centrosymmetric geometries; each box spans the full range measured over crystals and seeds for one arm of each matched core; parity forces the O(3) value to zero. c, Predicted magnitude on ten named compounds. Shaded rows are class m$\bar 3$m, where rotation alone forbids a response; on the two unshaded, non-cubic rows only parity protects the zero. Filled circle, O(3); open markers, the SO(3) arms. The axis is symmetric-log with a true zero, and every marker is plotted at its full-precision value from the run outputs.}
\label{fig:mechanism}
\end{figure}

\subsection*{The failure spans the released landscape}

The matched arms isolate the mechanism; an audit measures its prevalence. We classified eighteen released architectures by whether their features could, in principle, guarantee a symmetry-mandated zero for a parity-odd tensor, reading each source at a pinned version and probing numerically where inspection could not settle it. Six attach parity labels, so an odd-parity output vanishes structurally by construction. The remaining twelve lack the labels in three distinct ways (Table~\ref{tab:audit}): six are rotation-equivariant only, the class measured above; three carry a single unlabelled Cartesian vector channel, leaving the guarantee the practitioner's to build and unchecked; and three are invariant models. At random initialization the three deployed-generation models in the rotation-only class already show the rotation-ceiling signature, their m$\bar 3$m violation medians five to six orders of magnitude below the non-cubic ones. The three populations answer different questions, and pooling them understates the prevalence among the models the constraint can reach: of the twelve architectures that can produce a rank-3 tensor output as built, the six parity-aware and the six rotation-only, exactly half type parity. The full audit, with the deciding source line for each row, is in \snote{snote:audit}.

Two rows show why measurement, not reading, had to settle the classification. NequIP exposes a boolean named \texttt{parity}, yet disabling it does not remove parity labels, which is why our SO(3) arms relabel irreducible representations directly. GotenNet~\cite{aykent2025gotennet}, conversely, is not parity-typed yet passed the reflection test, moving from undetermined to parity-aware (Table~\ref{tab:audit} footnote).

\begin{table}[t]
\centering
\caption{Feature parity across eighteen released architectures. ICTP~\cite{zaverkin2024ictp} and GotenNet~\cite{aykent2025gotennet} were settled by numerical reflection test rather than by source, which did not decide them; eSEN, UMA and EqV3 were both source-read and reflection-measured. Groups follow the classification in the text; versions, commits and the deciding evidence for every row are tabulated in the Supplementary Information, and the classification does not extend to later releases.}
\label{tab:audit}
\fontsize{8bp}{10bp}\selectfont
\begin{tabular*}{\textwidth}{@{\extracolsep{\fill}}l l c@{}}
\toprule
Architecture & Internal features & Parity typed \\
\midrule
\multicolumn{3}{@{}l}{\textbf{Parity-aware} (6). An odd-parity output is structurally zero.} \\
NequIP~\cite{batzner2022nequip}      & O(3) irreps, natural parity $(-1)^l$ & yes \\
Allegro~\cite{musaelian2023allegro}  & O(3) irreps, natural parity $(-1)^l$ & yes \\
MACE~\cite{batatia2022mace}          & O(3) irreps, natural parity $(-1)^l$ & yes \\
MatTen~\cite{wen2024matten}          & O(3) irreps, natural parity $(-1)^l$ & yes \\
ICTP~\cite{zaverkin2024ictp}         & Cartesian harmonics; parity implicit & yes\textsuperscript{$\ast$} \\
GotenNet~\cite{aykent2025gotennet}   & untyped attention; polar vector output & yes\textsuperscript{$\ast$} \\
\addlinespace[3pt]
\multicolumn{3}{@{}l}{\textbf{Rotation-only} (6). Nothing can force the output to vanish.} \\
EquiformerV1~\cite{liao2023equiformer}   & SO(3) irreps, all declared even & no \\
EquiformerV2~\cite{liao2024equiformerv2}  & degree and order only; no parity flag & no \\
eSCN~\cite{passaro2023escn}               & degree and order only; no parity flag & no \\
eSEN~\cite{fu2025esen}                    & degree and order only; no parity flag & no\textsuperscript{$\ast$} \\
UMA~\cite{wood2025uma}                    & degree and order only; no parity flag & no\textsuperscript{$\ast$} \\
EquiformerV3~\cite{equiformerv3}          & degree and order only; no parity flag & no\textsuperscript{$\ast$} \\
\addlinespace[3pt]
\multicolumn{3}{@{}l}{\textbf{Vector-only} (3). Implicitly polar but untyped, and no rank-3 path.} \\
PaiNN~\cite{schutt2021painn}          & one Cartesian vector channel & untyped \\
TorchMD-Net~\cite{tholke2022torchmdnet} & one Cartesian vector channel & untyped \\
ViSNet~\cite{wang2024visnet}          & one Cartesian vector channel & untyped \\
\addlinespace[3pt]
\multicolumn{3}{@{}l}{\textbf{Invariant} (3). Cannot express an equivariant tensor output.} \\
SchNet~\cite{schutt2017schnet}            & interatomic distances only & n/a \\
DimeNet++~\cite{gasteiger2020dimenetpp}   & distances and bond angles & n/a \\
FAENet~\cite{duval2023faenet}             & not equivariant; frames applied externally & n/a \\
\bottomrule
\end{tabular*}

\vspace{2pt}
{\fontsize{7bp}{8.5bp}\selectfont
\textsuperscript{$\ast$}Confirmed by measurement at random initialization. GotenNet's vector output obeys the polar-vector law under improper operations to $1.8\times10^{-15}$, against $1.9$ for the pseudovector law; eSEN, UMA and EqV3 satisfy the rotation law to $10^{-5}$--$10^{-4}$ yet violate the mirror law by order unity. EquiformerV1 is classed by its released default configuration, although its formulation admits odd irreducible representations.\par}
\end{table}

The dedicated crystal-tensor predictors that enforce symmetry at the readout face the same test. Of the four current models, two release piezoelectric code and checkpoints, GMTNet~\cite{yan2024gmtnet} and CEITNet~\cite{ceitnet2026}, both built on an invariant backbone with equivariance synthesized only at the tensor head; we ran both on the same 2,000 crystals through their own pipelines (GoeCTP~\cite{goectp2024} and EATGNN~\cite{lu2025eatgnn} release only rank-2 dielectric implementations). Their head algebra does produce structural zeros for most crystals on exactly symmetric inputs, but both still false-flag a substantial minority with their space-group mask off, and they rely on that mask differently. GMTNet's mask repairs the idealized coordinates completely yet leaks on the raw ones, where the symmetry determination degrades. CEITNet's mask barely reduces its rate, being inert for the fully-forbidden point groups: it fails most on the m$\bar 3$m crystals every rotation-equivariant model here gets right by construction, traceable to a benchmark protocol~\cite{choudhary2020jarvis} that removes zero tensors from training. Readout-level repair works as far as its bookkeeping reaches; the guarantee either lives in the features or depends on an external oracle that real inputs erode (\snote{snote:sota}).

The decision also propagates to downstream property heads: as heads are increasingly attached to the frozen embeddings of pretrained universal potentials, the backbone's symmetry group is inherited by every target built on them. We trained a piezoelectric head on the frozen features of two released potentials under one protocol. On MACE-MP-0~\cite{batatia2023foundation}, whose features carry parity labels, the head's violations are exactly zero on every crystal and seed; on eSEN-30M-OAM~\cite{barrosoluque2024omat24}, whose features declare every degree even, the same construction false-flags at the full headline rate, at comparable regression accuracy. A parity-even rank-2 target such as the dielectric tensor never exposes this; the first parity-odd one does.
